\section{Compatible transfer and the one-charge ledger}

We now place activation, representation, and coherent reassembly on one
carrier. Fix a nonzero parent form \(R_{X,\nu}\), the common terminal map
\(L_X\) with \(Y_X\ne\{0\}\), a nonzero
represented map \(A_X:Y_X\to H_X\), and a missing/reassembly map
\(M_X:Y_X\to H_X\). Set
\[
 B_X=A_X+M_X.
\]
All maps and supports are fixed before \(z\).

Define the exact stage constants
\begin{align}
 \tau_{X,\nu}
 &:=\inf_{\mathcal P_{X,\nu}(z)>0}
       \frac{\|L_Xz\|^2}{\mathcal P_{X,\nu}(z)},
 \label{eq:tau-stage}\\
 \beta_X
 &:=\inf_{0\ne y\in Y_X}\frac{\|A_Xy\|^2}{\|y\|^2},
 \label{eq:beta-stage}\\
 \delta_X
 &:=\inf_{A_Xy\ne0}
       \frac{\|B_Xy\|^2}{\|A_Xy\|^2}.
 \label{eq:delta-stage}
\end{align}
The direct parent-to-final constant is
\begin{equation}\label{eq:direct-gamma}
 \Gamma_{X,\nu}
 :=\inf_{\mathcal P_{X,\nu}(z)>0}
 \frac{\|B_XL_Xz\|^2}{\mathcal P_{X,\nu}(z)}.
\end{equation}
The next product is the TPC-62 three-module mechanism specialized to the
common activated carrier; its role here is compatibility and loss provenance,
not a new abstract multiplication theorem.

\begin{theorem}[Compatible three-stage certificate]\label{thm:three-stage}
If \(\tau_{X,\nu},\beta_X,\delta_X\) are defined on the common carrier above,
then
\begin{equation}\label{eq:stage-product}
 \boxed{\quad
 \Gamma_{X,\nu}\ge\tau_{X,\nu}\beta_X\delta_X.
 \quad}
\end{equation}
The product can be sharp, but it need not equal the direct constant and can be
arbitrarily smaller.
\end{theorem}

\begin{proof}
If \(\mathcal P_{X,\nu}(z)>0\) and \(L_Xz=0\), then
\(\tau_{X,\nu}=0\) and the claim is trivial. Otherwise apply the three
definitions successively to the same vector \(z\), \(L_Xz\), and \(A_XL_Xz\).
If the latter vanishes, then either \(\beta_X=0\), or the ratio in
\eqref{eq:delta-stage} is outside its nonzero reference set; the displayed
bound again reduces to zero. Taking the infimum proves
\eqref{eq:stage-product}. One-dimensional scalar maps give equality.
\end{proof}

The exact reassembly constant \(\delta_X\) is the TPC-62 shorted generalized
eigenvalue relative to \(A_X^*A_X\). Its specialization to the current carrier
is recorded for compatibility; it is not a new proof of the master short.

\begin{proposition}[Shorting-collapse dichotomy]\label{prop:short-collapse}
If \(\beta_X>0\), then \(A_X\) is injective on \(Y_X\), so the null space in
the reassembly short is zero. Consequently
\begin{equation}\label{eq:ordinary-reassembly-gram}
 \delta_X=\lambda_{\min}\!\left[
 (A_X^*A_X)^{-1/2}B_X^*B_X(A_X^*A_X)^{-1/2}
 \right].
\end{equation}
If the short is genuinely nontrivial, then the uniform staged represented
gate has already failed, although the direct final constant may still be
positive.
\end{proposition}

\begin{proof}
By \cref{thm:row-gram}, \(\beta_X>0\) is equivalent to
\(\Ker A_X=\{0\}\) on \(Y_X\). The nuisance image used in the short is then
zero, so the projected Gram reduces to \eqref{eq:ordinary-reassembly-gram}.
Conversely a nonzero nuisance null space implies \(\beta_X=0\). Direct
positivity is logically independent: for example, on \(\mathbb C^2\), take
\(A(x_1,x_2)=(x_1,0)\), \(M(x_1,x_2)=(0,x_2)\); then \(A+M=I\).
\end{proof}

\subsection{Exponent bookkeeping}
Suppose, with uniform remainders on the same scope,
\[
 \tau_{X,\nu}\ge X^{-\lambda_{\rm act}+o(1)},\qquad
 \beta_X\ge X^{-\lambda_{\rm rep}+o(1)},\qquad
 \delta_X\ge X^{-\lambda_{\rm reasm}+o(1)}.
\]
Then
\begin{equation}\label{eq:exponent-ledger}
 \Gamma_{X,\nu}
 \ge X^{-(\lambda_{\rm act}+\lambda_{\rm rep}
                 +\lambda_{\rm reasm})+o(1)}.
\end{equation}
Parent inflation belongs entirely to \(\lambda_{\rm act}\). The represented
Gram is normalized by terminal energy and the reassembly Gram by represented
energy, so charging the same repair inflation in either later arrow would be a
double count. A downstream arithmetic endpoint still requires the complete
loss sum to be \emph{strictly} below its budget; equality at \(1/400\) is not a
certificate.

Nothing in this paper proves any of the three displayed polynomial bounds for
the physical unlabelled fixed-shift map. Formula \eqref{eq:exponent-ledger} is
a compatibility theorem, not a claimed endpoint estimate.
